\scanpage{132}
\uline{Corollaire} Fonctorialité \figure{\begin{tikzpicture}\node[draw,ellipse,inner sep=5pt] {en $X_K$};\end{tikzpicture}} de l'ensemble $\mathrm{Ép}\deleted{\unsure{\mathrm{ais}}}(\bar X)$ des points \deleted{\unsure{de}} épais de $\bar X$ [\unsure{munis} de leurs corps résiduels $\ldots$], pour des \deleted{morphismes} \inserted{applic. rationnelles} $X_K-Y_K$ \uline{dominantes gén.\textsuperscript{t} séparables de} $K$-\uline{schémas} \uline{propres}.

\uline{Proposition 2} Soit $S$ une partie constructible de $\bar X$, $X_K$-simple, [tel que $\unsure{\ell|d}(\unsure{W|V},x)$ y soit borné], et tel que $\bar\varphi(S)$ existe [p.ex $Y$ propre sur $V$] et soit $Y_K$-simple.

Alors $\bar\varphi(S)$ est constructible. De plus, il y a une partition \inserted{finie} $S=\coprod_i S_i$, les $S_i$ des sous-schémas irréductibles constructibles de $X$, telle que $\varphi$ soit morphique d'indice fini sur chaque $S_i$, [et $\varphi(S_i)$ est alors un ens. constructible irréductible $\ldots$]

\uline{Corollaire} Fonctorialité en $X_K$ variant comme dessus, et de plus simple sur $K$, de l'ens. $\mathrm{Cons}(\bar X)$ des parties constructibles de $\bar X$.
\margin{131}
